\documentclass[11pt]{article}
\usepackage[margin=1in]{geometry}
\usepackage{amsmath,amssymb,amsthm,hyperref}
\newtheorem{lemma}{Lemma}
\newtheorem{theorem}[lemma]{Theorem}
\title{The ordinary gap-six terminal family $N=7^r+2$}
\author{Internal research note}
\date{18 November 2026}
\begin{document}\maketitle
Fix $r\ge1$, $L=7^r$, $N=L+2$, and $c=N+5$.
Throughout $N$ denotes the terminal order.
Let $K=\mathbb C_7$, with valuation ring $\mathcal O$ and residue field
$k=\mathcal O/\mathfrak m$.  Put $A=c-zD_z+z$ and
$\Pi_N(X)=\prod_{\ell=0}^{N-1}(X+6\ell)$.
For $Q\in\mathcal O[X]$, $\deg Q\le4$, define
\[
 G_Q(X)=\binom{X/6+N-1}{N}Q(X),\qquad
 d_j(Q)=\sum_{\ell=0}^j(-1)^\ell\binom j\ell G_Q(c-\ell),
\]
and the six boundary rows at $z^6=7v$ by
\[
 V_a(Q;v)=\sum_{\ell\ge0}\frac{7^\ell a!}{(a+6\ell)!}
                   d_{a+6\ell}(Q)v^\ell,\qquad0\le a<6.       \tag{1}
\]
All sums are finite. The binomial polynomial is integer-valued on
$\mathbb Z_7$, by continuity from nonnegative integers.
Its arguments at $c-\ell$ lie in $\mathbb Z_7$.
Consequently $d_j(Q)\in\mathcal O$ and $V_a(Q)\in\mathcal O[v]$:
the coefficient weight in (1) has valuation
$(s_7(a+6\ell)-a)/6\ge0$.
The basis-column matrix belongs to $\mathbb Z_7[v]$.
The finite formal identity
\[
 \frac{\Pi_N(A)Q(A)1}{6^NN!}=\sum_jd_j(Q)\frac{z^j}{j!}
\]
shows that simultaneous vanishing of the six rows is precisely the
corresponding cyclic rank relation.

\begin{lemma}[Uniform seeds]\label{lem:seeds}
Put
\[
 S=S_r(v)=1+\sum_{i=1}^r(-1)^i v^{(7^i-1)/6}.
\]
The rows $V_0,V_1$ are divisible by7 on all five basis columns.
For $R_a=\overline{V_a/7}$, $a=0,1$, the six seeds are
\[
 R_0(1)=1,\quad R_1(1)=5+4S,\quad
 \bar V_2(1)=1+2S,\quad \bar V_3(1)=4S,\quad
                         \bar V_4(1)=\bar V_5(1)=0.           \tag{2}
\]
\end{lemma}
\begin{proof}
Use the exact diagonal generating function from
\texttt{lem:diagonal} of
\href{GAP_SIX_TWO_POWER_PLUS_ONE_SFC.tex}
{the second thin-family note}:
\[
 \sum_{n\ge0}\frac{B_{n,n,6}(z)}{6^nn!}t^n
   =\kappa\exp(z(1-R)),\quad
 R-R^7=t,\ R(0)=1,\quad\kappa=\frac{6R}{7R^6-1}.
                                                               \tag{3}
\]
Here $B_{n,\gamma,6}=\Pi_n(\gamma+5-zD_z+z)1$ and
$\kappa\equiv R-7t\pmod{49}$.

For completeness, we spell out the coefficient extraction needed
at this new terminal family. All following equalities are formal
$t$-adic identities, never evaluations of an exponential at the
boundary. Set
\[
 T=\sum_{i\ge0}t^{7^i},\quad
 C(T)=\frac{T^7-\sum_{i\ge0}t^{7^{i+1}}}{7},\quad
 P(T)=T+3T^2+5T^3+5T^4+3T^5+T^6.
\]
Then $R=1+T+7R_1$ and $R_1\equiv C(T)+P(T)\pmod7$.
Define integer polynomials
$a_0=-1$, $a_i=-1-va_{i-1}^7$, and
$A_0=\sum_i a_it^{7^i}$.
Thus $\bar a_r=-S_r(v)$.
The integral Artin--Hasse series, imported as
\texttt{lem:artin-hasse-closure} of the
\href{ARITHMETIC_FRONTIER_TERMINAL_ORDERS.tex}
{arithmetic note}, gives a unique integral formal series $A_*$ with
\[
 A_*+vA_*^7+7^6v^8A_*^{49}+\cdots=1-R,\qquad
 \exp(z(1-R))=\mathcal E_7(zA_*)\quad(z^6=7v).
\]
Writing $A_*=A_0+7A_1$, one obtains
\[
 A_1\equiv-R_1-vC(A_0),\qquad
 C(A_0)=\frac{A_0^7-\sum_i a_i^7t^{7^{i+1}}}{7}.
\]
The cyclic sections of the exponential satisfy
\[
 \mathcal C_{6,0}\exp\equiv1+42vA_0^6\pmod{49},\qquad
 \mathcal C_{6,1}\exp\equiv A_*+14vA_*^7\pmod{49};
\]
the second uses $[X^7]\mathcal E_7=103/720\equiv2\pmod7$.
For $0\le a<6$, reduction of $a!\mathcal C_{6,a}\exp$ is $A_0^a$.

The integer $N=L+2$ has base-seven digit sum3.
It is not a sum of at most two powers of7, so
$[t^N](1+T)=[t^N](1+T)A_0=0$ exactly. Hence
\[
 \begin{split}
 R_0(1)&=[t^N]\{R_1-t+6v(1+T)A_0^6\},\\
 R_1(1)&=[t^N]\{R_1(A_0-1-T)-v(1+T)C(A_0)
                           +2v(1+T)A_0^7-tA_0\}.
 \end{split}                                                   \tag{4}
\]
All equations in (4) are in characteristic7.
The carry series $C(T),C(A_0)$ are supported at exponents1 modulo6.
Their products here have exponents1 or2 modulo6, whereas $N=3$
modulo6. Terms with $A_0^7$ reduce to sums of at most two powers
of7 and also vanish. The $A_0^6$ terms have exponents0 or1 modulo6.
Thus the first line of (4) is $[t^N]5T^3=15=1$.
Only $-5T^3+3T^2(A_0-T)$ survives in the second line.
There is no carry in selecting three powers of7, and
\[
 [t^N]T^3=3,\qquad [t^N]T^2A_0=a_r+2a_0=a_r-2.
\]
The result is $3a_r-30=5+4S$ in characteristic7.

The ordinary seed formula is
$\bar V_a(1)=[t^N](1+T)A_0^a$.
For $a=2$ the term $TA_0^2$ gives $1-2a_r=1+2S$;
for $a=3$ the term $A_0^3$ gives $3a_r=4S$;
the remaining seeds vanish by their number of selected powers.

Finally the exact raising recurrence is
\[
 \begin{split}
 V_0(XQ)&=(c-6v\partial_v)V_0(Q)+7vV_5(Q)/5!,\\
 V_a(XQ)&=(c-a-6v\partial_v)V_a(Q)+aV_{a-1}(Q),\quad1\le a<6.
 \end{split}                                                   \tag{5}
\]
Starting from the two divisible seeds, (5) proves global
7-divisibility for $V_0,V_1$ on every required basis column.
\end{proof}

\begin{lemma}[A fixed universal matrix]\label{lem:matrix}
In the column basis $1,X,\ldots,X^4$, the six reduced rows
$R_0,R_1,\bar V_2,\ldots,\bar V_5$ form
\[
 M(S,v)=
 \begin{pmatrix}
 1&0&0&3Sv&Sv-v\\
 -3S-2&-1&1&-1&3Sv+1\\
 2S+1&3-2S&2S+3&-2S-1&2S-3\\
 -3S&-2S-1&-2S&3S-3&2S+3\\
 0&2S&1&3-S&S-2\\
 0&0&3S&2S+2&S+3
 \end{pmatrix}.                                                \tag{6}
\]
\end{lemma}
\begin{proof}
One has $vS_r'(v)=S_r(v)-1$ in characteristic7.
Apply (5) to (2), using
$\delta=v\partial_v+(S-1)\partial_S$ and $c=0$ in the residue field.
The recurrences for the six rows are respectively
\[
 R_0(XQ)=\delta R_0(Q)+v\bar V_5(Q),\quad
 R_1(XQ)=(-1+\delta)R_1(Q)+R_0(Q),
\]
\[
 \bar V_2(XQ)=(-2+\delta)\bar V_2(Q),\quad
 \bar V_a(XQ)=(-a+\delta)\bar V_a(Q)+a\bar V_{a-1}(Q)
                                                  \quad(3\le a<6).
\]
These generate exactly (6).
\end{proof}

\begin{lemma}[A complete six-state obstruction]\label{lem:certificate}
For every $r\ge1$ and every $v\in\bar{\mathbb F}_7$,
the matrix $M(S_r(v),v)$ has rank5.
\end{lemma}
\begin{proof}
Let $m_i$ be the determinant obtained from (6) by omitting row $i$,
numbered 0 through 5. Put
\[
 \begin{split}
 W(S)&=S^6+S^5+2S^4+3S^3+4S^2+5S+5,\\
 q(S)&=S^{10}+3S^9+4S^6+6S^5+6S^4+6S^3+5S^2+2S+3.
 \end{split}
\]
Exact polynomial identities in $\mathbb F_7[S,v]$ are
\[
 \begin{split}
 v-W&=S m_0+(6S+4)m_1,\\
 q&=(6S^5+5S^4+4S^3+3S^2+2S+1)m_0\\
   &\hspace{1cm}+(S^5+5S^4+2S^3+6S^2+3S)m_1.
 \end{split}                                                   \tag{7}
\]
Hence rank drop forces $q(S)=0$ and $v=W(S)$.

In $\mathbb F_7[S]/(q)$ one verifies $W^8+W+1=0$.
Consequently the recurrence $x_0=1$,
$x_{j+1}=1-Wx_j^7$ has the complete cycle
\[
 1,\quad1-W,\quad-2W,\quad-1-2W,\quad-1-W,\quad0,\quad1.          \tag{8}
\]
Euclid's algorithm verifies
$\gcd(q,S-x_j)=1$ for all six states in (8).
The literal matrix, both six-minor multiplier lists, and the complete
six-state orbit are retained in
\texttt{gap\_six\_power\_plus\_two\_}\allowbreak\texttt{certificate.json}.
The accompanying checker reconstructs all 30 matrix entries from
the seeds, computes all six maximal minors, and checks both
membership identities (7), every state transition, and the cycle return.
For an actual root of $q$ and $v=W(S)$, the sequence $S_r(v)$
is precisely the specialization of this recurrence.
It can therefore never equal $S$. This contradicts rank drop.
The quotient need not be a field; neither its irreducibility nor
an unproved finite-orbit bound is used.
\end{proof}

\begin{theorem}[The terminal family]\label{thm:family}
Every nonzero ordinary univariate binomial $X(A+BX^6)$ satisfies the
consecutive-moment coprimality test at every terminal order
$N=7^r+2$, $r\ge1$.
Equivalently, when $AB\ne0$, the normalized factorial-moment
polynomials at orders $N-1,N$ have no common complex root.
\end{theorem}
\begin{proof}
Degenerate monomials satisfy the assertion directly; assume $AB\ne0$.
The standard phase normalization is then nonzero.
Write
$T(t)=\sum_{j=0}^N\binom Nj(N)_{6j}t^j$.
For the ordinary functional $\mathcal L(X^j)=j!$, and
$\theta=tD_t$, the moment polynomials after putting $t=B/A$ are
\[
 \begin{split}
 \mathcal L(f^{N-1})&=A^{N-1}(N-1)!\,(T-\theta T/N),\\
 \mathcal L(f^N)&=A^NN!\,(T+6\theta T/N).
 \end{split}
\]
Indeed $\binom{N-1}{j}=(1-j/N)\binom Nj$ and
$(N+1)_{6j}=(1+6j/N)(N)_{6j}$.
The two displayed combinations have determinant $7/N$ relative
to $T,\theta T$.
A common root is thus exactly a multiple root of$T$.
By \texttt{thm:radius} of
\href{PHYSICAL_PRIME_TOTAL_DEGREE_REDUCTION.tex}
{the physical prime-total-degree note},
every hypothetical multiple root has $u=-1/t$ of valuation1.
Thus $v=u/7$ is a unit.

The necessary cyclic rank relation is
\texttt{lem:critical-ideal} of
\href{gap_three_hurwitz_structure.tex}
{the adjoint-structure note}, specialized to gap6 and
the fixed physical shape$\gamma=N$.
Although its exponential proof is over$\mathbb C$, the cyclic
minors here are rational polynomials. Any multiple root is
algebraic; embed it into$\mathbb C$, apply that necessary rank
relation, and transport the resulting rational polynomial
identities to its chosen embedding into$\mathbb C_7$.
No exponential is evaluated at the 7-adic boundary.
The polynomial basis used in (1) agrees with the consecutive
moment basis by an integral triangular change: for $0\le j\le4$,
\[
 B_{N+j,N,6}
 =\Pi_N(A)\prod_{\ell=N}^{N+j-1}(A+6\ell)1 .
\]
The empty product is 1 and each product is monic of degree $j$.
The five relation coefficients can therefore be normalized to a
primitive vector in$\mathcal O^5$.
All six rows in (1) then vanish; division of the globally divisible
first two rows preserves their vanishing.
Reduction would give a nonzero kernel vector of
$M(S_r(\bar v),\bar v)$.
The preceding lemma excludes this.
\end{proof}

The theorem is an infinite terminal family, not all-index gap-six
SFC. The remaining ordinary gap-six orders, unrestricted binomial
SFC, and the full factorial conjectures remain open here.
The checker verifies polynomial identities and complete finite
certificates; it is not formal proof verification or external review.
\end{document}
